\documentclass[12pt,a4paper]{article}

% Компиляция: pdflatex matematicheskaya_postanovka_car_damage.tex
\usepackage[T2A]{fontenc}
\usepackage[utf8]{inputenc}
\usepackage[russian]{babel}
\usepackage[a4paper,left=25mm,right=20mm,top=20mm,bottom=20mm]{geometry}
\usepackage{amsmath,amssymb,amsfonts,mathtools}
\usepackage{microtype}
\usepackage{enumitem}
\usepackage{indentfirst}
\usepackage{hyperref}
\hypersetup{hidelinks}
\setlength{\parindent}{1.25cm}
\setlength{\parskip}{0pt}
\linespread{1.15}

\newcommand{\E}{\mathbb{E}}
\newcommand{\Pbb}{\mathbb{P}}
\newcommand{\R}{\mathcal{R}}
\newcommand{\Loss}{\mathcal{L}}
\DeclareMathOperator{\Var}{Var}

\begin{document}

\section{Постановка задачи}


Изображением автомобиля считаем RGB-изображение, хранящееся 
в в виде трёхмерного массива $X\in[0,1]^{H\times W\times 3}$, 
где $H$ --- высота, $W$ --- ширина, а три канала 
соответствуют цветам R, G и B. На изображении могут 
присутствовать повреждения автомобиля, которые 
необходимо локализовать и классифицировать.
Под локализацией будем иметь в виду ограничивающий область повреждения прямоугоольник. 
Он описывается набором
 $\{x, y, w, h\}$, где $(x, y)$ 
--- координаты верхнего левого угла ограничивающего 
прямоугольника, $w$ --- ширина, $h$ --- высота.
Каждый ограничивывающий прямоугольник соответствует 
одному из классов повреждений, 
уверенность модели в данном классе характеризуется числом из интервала $[0,1]$.


Тогда пусть задана обучающая выборка изображений автомобилей
\begin{equation*}
    \mathcal D_{\mathrm{train}} = \{(X_i,Y_i)\}_{i=1}^{N},
\end{equation*}
где $N$ --- количество изображений; $X_i\in[0,1]^{H\times W\times 3}$ --- изображение;
$Y_i$ --- множество аннотаций повреждений на $i$-м изображении.

$Y_i$ имеет следующую структуру: 
\begin{equation*}
    Y_i = \{y_{ij}\}_{j=1}^{n_i},
    \qquad
    \tilde{y}_{ij} = (b_{ij},c_{ij},p_{ij},m_{ij},s_{ij}),
\end{equation*}
где $n_i$ --- число повреждений, а
 компоненты $\tilde{y}_{ij}$ описываются:
\begin{itemize}[leftmargin=*,itemsep=2pt,topsep=3pt]
    \item $b_{ij}=(x_{ij},y_{ij},w_{ij},h_{ij})$ --- ограничивающий прямоугольник повреждения;
    \item $c_{ij}\in\{1,\ldots,K\}$ --- класс повреждения;
    % \item $p_{ij}\in\{1,\ldots,P\}$ --- повреждённая деталь автомобиля;
    % \item $m_{ij}\in\{0,1\}^{H\times W}$ --- бинарная маска повреждения;
    % \item $s_{ij}\in\{1,\ldots,S\}$ --- категория тяжести повреждения.
\end{itemize}


Требуется построить параметризованную модель
\begin{equation*}
    f_\theta:\mathcal X\to\mathcal Y,
\end{equation*}
которая по изображению $X \in \mathcal X$, где $X \in[0,1]^{H\times W\times 3}$;
$\mathcal X$ - множество всех возможножных изображений, 
предсказывает множество повреждений
\begin{equation*}
    \widehat Y=f_\theta(X)
    =\{(\widehat b_j,\widehat c_j,\widehat q_j)\}_{j=1}^{\widehat n},
\end{equation*}
где $\widehat b_j$ --- предсказанная область повреждения, $\widehat c_j$ --- его тип, $\widehat q_j\in[0,1]$ --- оценка уверенности детектора.

Параметры $\theta$ определяются минимизацией функции потерть $\Loss_{\mathrm{det}}$:
\begin{equation*}
    \theta^*=\underset{\theta}{\arg\min}\;
    \E
    \left[\Loss_{\mathrm{det}}(f_\theta(X),Y) |  X = X_{train}\right].
    \label{eq:det-main}
\end{equation*}

В общем случае функцию потерь детектора можно записать как
\begin{equation}
    \Loss_{\mathrm{det}}
    =\lambda_1\Loss_{\mathrm{cls}}
    +\lambda_2\Loss_{\mathrm{box}}
    % +\lambda_3\Loss_{\mathrm{aux}},
\end{equation}
где $\Loss_{\mathrm{cls}}$ отвечает за классификацию
 повреждений, $\Loss_{\mathrm{box}}$ --- за локализацию.
%  а $\Loss_{\mathrm{aux}}$ содержит дополнительные 
%  архитектурно-зависимые компоненты (например, потери для дискретизированной регрессии границ). 
Коэффициенты $\lambda_1,\lambda_2 \geq 0 $ задают относительный вклад слагаемых.

% В качестве функции потерь классификации будем использовать
% бинарную кросс-энтропию:
% \begin{equation*}
%     \mathcal{L}_{\mathrm{BCE}}
%     =
%     -\frac{1}{RK}
%     \sum_{r=1}^{R}
%     \sum_{k=1}^{K}
%     \left[
%         t_{rk}\ln\hat{p}_{rk}
%         +
%         (1-t_{rk})\ln(1-\hat{p}_{rk})
%     \right].
% \end{equation*}

% Здесь для фиксированного изображения $X_i$:
%     $R$ --- количество кандидатов в повреждения,
%     формируемых нейросетью до постобработки;
%      $K$ --- количество классов повреждений;
%      $\hat{p}_{rk}\in(0,1)$ --- предсказанная
%     вероятность принадлежности $r$-го кандидата
%     к $k$-му классу повреждений;
%      $t_{rk}\in\{0,1\}$ --- целевое значение
%     для $r$-го кандидата и $k$-го класса,
%     определяемое при сопоставлении предсказаний
%     с аннотациями $Y_i$.

% Минимизация $\mathcal{L}_{\mathrm{BCE}}$ направлена
% на повышение точности классификации
% обнаруженных повреждений.

% Для оценки точности локализации повреждений
% рассмотрим метрику Intersection over Union (IoU):
% \begin{equation*}
%     \operatorname{IoU}(b_{ij},\hat{b}_r)
%     =
%     \frac{|b_{ij}\cap\hat{b}_r|}
%          {|b_{ij}\cup\hat{b}_r|}.
% \end{equation*}

% Здесь $b_{ij}$ --- истинный ограничивающий
% прямоугольник $j$-го повреждения на изображении
% $X_i$, а $\hat{b}_r$ --- прямоугольник,
% предсказанный нейросетью для $r$-го кандидата.
% Величины $|b_{ij}\cap\hat{b}_r|$
% и $|b_{ij}\cup\hat{b}_r|$ обозначают площади
% пересечения и объединения прямоугольников
% соответственно.

% Метрика IoU принимает значения из отрезка
% $[0,1]$. Значение $1$ соответствует полному
% совпадению прямоугольников, а значение $0$
% означает отсутствие их пересечения.

% Соответствующая функция потерь имеет вид:
% \begin{equation*}
%     \mathcal{L}_{\mathrm{IoU}}
%     (b_{ij},\hat{b}_r)
%     =
%     1-\operatorname{IoU}(b_{ij},\hat{b}_r).
% \end{equation*}

% Недостатком данной функции является отсутствие
% информации о расстоянии между прямоугольниками
% в случае, когда они не пересекаются.

% Поэтому для обучения модели будем использовать
% Complete Intersection over Union (CIoU),
% учитывающую степень пересечения прямоугольников,
% расстояние между их центрами и различие
% соотношений сторон.

% Функция потерь CIoU определяется выражением:
% \begin{equation*}
%     \mathcal{L}_{\mathrm{CIoU}}
%     (b_{ij},\hat{b}_r)
%     =
%     1-\operatorname{IoU}(b_{ij},\hat{b}_r)
%     +
%     \frac{\rho^2(z_{ij},\hat{z}_r)}{d^2}
%     +
%     \alpha v,
% \end{equation*}

% где
%     $z_{ij}$ и $\hat{z}_r$ --- центры истинного
%     и предсказанного прямоугольников соответственно;
%     $\rho(z_{ij},\hat{z}_r)$ --- евклидово
%     расстояние между центрами прямоугольников;
%     $d$ --- длина диагонали минимального
%     прямоугольника, содержащего оба рассматриваемых
%     прямоугольника;
%     $v$ --- величина, характеризующая различие
%     соотношений сторон прямоугольников;
%     $\alpha$ --- коэффициент, определяющий
%     влияние различия соотношений сторон.


% Поскольку ранее координаты прямоугольника
% были определены как координаты его верхнего
% левого угла и размеры, координаты центров имеют вид:
% \begin{equation}
%     z_{ij}
%     =
%     \left(
%         x_{ij}+\frac{w_{ij}}{2},
%         y_{ij}+\frac{h_{ij}}{2}
%     \right),
% \end{equation}

% \begin{equation}
%     \hat{z}_r
%     =
%     \left(
%         \hat{x}_r+\frac{\hat{w}_r}{2},
%         \hat{y}_r+\frac{\hat{h}_r}{2}
%     \right).
% \end{equation}

% Величины $v$ и $\alpha$ вычисляются по формулам:
% \begin{equation}
%     v=
%     \frac{4}{\pi^2}
%     \left(
%         \arctan\frac{w_{ij}}{h_{ij}}
%         -
%         \arctan\frac{\hat{w}_r}{\hat{h}_r}
%     \right)^2,
% \end{equation}

% \begin{equation}
%     \alpha=
%     \frac{v}
%     {1-\operatorname{IoU}(b_{ij},\hat{b}_r)+v}.
% \end{equation}

% При полном совпадении прямоугольников
% полагаем $\alpha=0$.

% Пусть $\mathcal{P}_i$ --- множество индексов
% положительных кандидатов для изображения $X_i$,
% а $a_i(r)$ --- индекс истинного повреждения,
% сопоставленного кандидату $r$.

% Тогда функция потерь локализации
% на изображении $X_i$ определяется как:
% \begin{equation}
%     \mathcal{L}_{\mathrm{box}}
%     =
%     \frac{1}{Z_{\mathrm{box}}}
%     \sum_{r\in\mathcal{P}_i}
%     \omega_r
%     \mathcal{L}_{\mathrm{CIoU}}
%     \left(
%         b_{i,a_i(r)},\hat{b}_r
%     \right),
% \end{equation}

% где $\omega_r\geq0$ --- вес положительного
% кандидата, а $Z_{\mathrm{box}}>0$ ---
% нормировочный коэффициент.



% Для оценки качества будет использоваться средняя точность $\mathrm{mAP}@[0.5:0.95]$.

\section{Обзор литературы}
\subsection{Enhancing Used Car Price Predictions with Machine Learning-Based Damage Detection}


В работе R.~V.~Kulkarni и др.
предложена система автоматической оценки стоимости подержанных
автомобилей с учетом повреждений, обнаруженных на фотографиях.
Разработанная система состоит из двух частей: детекции
повреждений и прогнозирования стоимости автомобиля.

Для локализации повреждений используется  YOLOv5,
обученная на наборе из 4000 изображений.
Рассматриваются четыре класса повреждений: царапины
, деформации, разрушения деталей
 и повреждения стекол.
Помимо определения класса и координат ограничивающего
прямоугольника, система формирует показатель тяжести
повреждения $severity \in [0.1; 0.99] $.

Для прогнозирования стоимости автомобиля сравнивались
линейная регрессия, XGBoost и AdaBoost. Обучение моделей
проводится на наборе данных из 8000 автомобилей.
Наилучшие результаты продемонстрировал алгоритм AdaBoost,
для которого получены значения
$R^2 = 0.93$ и $\mathrm{MAPE} = 16.96\%$.

Таким образом, работа демонстрирует возможность совместного
использования методов детекции объектов и регрессионного
моделирования для оценки стоимости поврежденных автомобилей.

\subsection{YOLOv8, YOLOv9, and YOLOv10: A Study in
Automated Vehicle Damage Detection }
В работе S.~Dolhpolov и др.
проведено сравнительное исследование архитектур YOLOv8,
YOLOv9 и YOLOv10 для автоматического обнаружения
повреждений автомобилей. Основной целью исследования
является оценка точности и вычислительной эффективности
нейросетевых детекторов в задачах визуального контроля.

Для обучения использован открытый набор данных
Car Dents, содержащий 7258 изображений автомобилей.
Рассматриваются три класса повреждений: вмятины,
царапины и повреждения вследствие ДТП.
Обучение моделей проводилось с использованием
предобученных на открытых данных Common Objects in Context (COCO) весов и transfer learning. 

Экспериментальное сравнение показало, что наилучшие
результаты достигаются архитектурой YOLOv10:
$\mathrm{mAP@0.5} = 0.6508$,
$\mathrm{mAP@0.5:0.95} = 0.3492$.
Для YOLOv8 и YOLOv9 значения $\mathrm{mAP@0.5}$
составили 0.5969 и 0.6297 соответственно.
При этом среднее время обработки изображения
для YOLOv10 составило 14.1 мс против 12.5 мс
для YOLOv8. 

Дополнительно показана эффективность переноса обучения:
использование предобученных весов позволило повысить
$\mathrm{mAP@0.5}$ относительно обучения с нуля
на 22.5\%, 20.4\% и 16.9\% для YOLOv8,
YOLOv9 и YOLOv10 соответственно.

Таким образом, исследование демонстрирует преимущества
YOLOv10 при детекции повреждений автомобилей
и подтверждает целесообразность применения переноса
обучения. Полученные результаты могут использоваться
при выборе базовой архитектуры детектора для последующего
исследования методов локализации и классификации
повреждений автомобилей.

\subsection{Advanced Applications of Generative AI in
Actuarial Science: Case Studies Beyond ChatGPT}
В работе S.~Hatzesberger и I.~Nonneman
исследуется применение генеративного искусственного интеллекта
в страховой деятельности, в частности использование мультимодальных
больших языковых моделей (Vision-Language Models) для автоматической
классификации повреждений автомобилей.

Авторы сравнивают три подхода: сверточную нейронную сеть (CNN),
предобученную модель GPT-4o и модель GPT-4o, дополнительно обученную
на специализированном наборе изображений.

Эксперименты проводятся на выборке из 1500 изображений,
содержащих шесть классов повреждений: трещины ,
царапины , спущенные шины , вмятины,
разбитые стекла и поврежденные фары. 

По результатам экспериментов дообученная GPT-4o достигла
$\mathrm{Accuracy} = 0.880$ и
$\mathrm{Weighted}\ F_1 = 0.880$.
Для исходной GPT-4o значения составили 0.823 и 0.825,
а для CNN --- 0.837 и 0.835 соответственно.
Таким образом, дообучение позволило повысить качество
классификации мультимодальной модели.

Дополнительно исследуется способность GPT-4o определять
локализацию повреждения в текстовом формате, например
указывать наличие вмятины на заднем бампере или
разбитого лобового стекла.

Таким образом, работа демонстрирует перспективность применения
мультимодальных языковых моделей для анализа автомобильных
повреждений. Однако для полноценного решения задачи детекции
необходимо дополнительно исследовать возможность точной
пространственной локализации поврежденных областей и сравнить
данный подход со специализированными детекторами объектов.

\subsection{YUNet\_LLMClaimReport: An Enhanced
Automobile Insurance Fraud Detection and
Automated Claim Report Generation Using LLM}
В работе P.~Anand Kumar и S.~Sountharrajan
предложена гибридная система
YUNet\_LLMClaimReport для автоматизации обработки
страховых случаев, включающая детекцию и сегментацию
повреждений автомобилей, оценку страховых выплат
и генерацию отчетов.

Для обнаружения повреждений используется архитектура
YOLOv11, а для пиксельной сегментации поврежденных областей
применяется U-Net. Совместное обучение моделей осуществляется
с использованием комбинированной функции потерь:
\[
    \mathcal{L}_{total}
    = \lambda_{obj}\mathcal{L}_{YOLO}
    + \lambda_{seg}\mathcal{L}_{UNet},
\]
где $\lambda_{obj}$ и $\lambda_{seg}$ определяют
относительный вклад задач детекции и сегментации.

Обучение проводилось на собственном наборе данных,
содержащем 4000 изображений поврежденных автомобилей
с разметкой ограничивающих прямоугольников и масок
сегментации. Для оценки суммы страхового возмещения
используется линейная регрессия, учитывающая площадь
повреждений, тип автомобиля, тяжесть аварии
и географическое расположение.

Дополнительно применяется дообученная языковая модель
GPT-3.5-turbo, предназначенная для автоматического
формирования страховых отчетов на основе результатов
анализа изображений и текстовых описаний происшествий.

Предложенная архитектура достигла
$\mathrm{mAP@0.5}=0{,}9424$ и $\mathrm{IoU}=0{,}85$.
При прогнозировании страховых выплат получены
$\mathrm{MAE}=1250$ египетских фунтов и
$\mathrm{RMSE}=1800$ египетских фунтов.
Качество генерации отчетов характеризуется значениями
$\mathrm{BLEU}=0{,}85$.

Таким образом, исследование демонстрирует перспективность
совместного применения детекции и сегментации для оценки
автомобильных повреждений. Особый интерес представляет
использование площади сегментированных областей
в качестве признака модели прогнозирования стоимости
страхового возмещения.
\end{document}